\worksheettitle{M07\(\star\). Считаем без списка}{\modulemark{M07\(\star\)} Усложнение}
\studentnameblock

\begin{notebox}[От строк таблицы к количеству]
  Для условия $b=a+k$ цифра десятков $a$ начинается с 1. Чтобы $b$ не
  превысила 9, должно выполняться $a+k\leq9$, то есть $a\leq9-k$.
\end{notebox}

\begin{practicebox}[Разберите закономерность]
  \begin{center}
  \renewcommand{\arraystretch}{1.5}
  \begin{tabular}{c|c|c}
    \toprule
    условие & допустимые $a$ & число вариантов \\
    \midrule
    $b=a+1$ & от 1 до 8 & 8 \\
    $b=a+2$ & от 1 до \answerblank[10mm] & \answerblank[12mm] \\
    $b=a+4$ & от 1 до \answerblank[10mm] & \answerblank[12mm] \\
    $b=a+7$ & от 1 до \answerblank[10mm] & \answerblank[12mm] \\
    \bottomrule
  \end{tabular}
  \end{center}
  Сформулируйте правило количества вариантов через $k$:
  \answerblank[90mm].
\end{practicebox}

\begin{practicebox}[Сумма цифр]
  Для $a+b=S$ и $S\leq9$ цифра $a$ может принимать значения от 1 до
  $S$, а $b=S-a$ остаётся цифрой.

  Не выписывая числа, найдите количество вариантов:
  \begin{enumerate}[label=\textbf{\arabic*.}]
    \item сумма 4: \answerblank[15mm];
    \item сумма 7: \answerblank[15mm];
    \item сумма 9: \answerblank[15mm].
  \end{enumerate}
  Почему количество совпадает с суммой $S$? \answerblank[60mm].
\end{practicebox}

\begin{warningbox}[Граница правила]
  Проверьте сумму 10. Можно ли просто сказать, что вариантов 10? Учтите, что
  $b$ не может быть 10, а $a$ не может быть 0. Сделайте вывод:
  \answerblank[70mm].
\end{warningbox}

\begin{checkpointbox}[Связь с перебором]
  Объясните, почему подсчёт допустимых значений $a$ доказывает полноту так же,
  как заполненная таблица: \answerblank[65mm].
\end{checkpointbox}
